\begin{remark}[Limits of automatic analytic comparison]
The abstract identity $\Aut(X)=\C^*$ does not supply a
meromorphic compactification of the action. On an ordinary
two-torus the distinguished orbit is a nonconstant holomorphic
map from $\C^*$. If it extended meromorphically across both
ends to $\PP^1$, it would extend holomorphically: a
meromorphic map from a smooth curve to a compact complex
manifold has a removable indeterminacy. Every holomorphic
map $\PP^1\to T^2$ is constant, since it lifts to $\C^2$.
This is a contradiction. Thus an algebraic structure on the
abstract automorphism group is insufficient here.

The compactification theorem of
\citet[Theorem 1.2]{BiliottiGhigi2021Root} assumes Fujiki
class $\mathcal C$; that hypothesis is unavailable.
The local isomorphism theorem of
\citet[Theorem 3]{Meersseman2011Root} assumes pointwise
isomorphism on the whole base, including its central point,
and constant vector-field dimension. It cannot identify the
unknown central fibre. The Hopf compactification example of
\citet[Example 1]{KatzarkovEtAl2025Root} likewise illustrates
that an analytic open part need not determine a bimeromorphic
compactification. Their definability-based comparison is not
used: no definability of our punctured comparison has been
proved. B.4.7 supplies a proper family exhibiting the relevant
failure directly. None of these comparisons closes B.4.

Non-K\"ahler modifications of elliptic degenerations are
established. The product case in
\citet[Section 3, Example 3.4(1)]{Fujimoto1988Root}, using
the preceding construction, already gives such degenerations
with elliptic abelian nearby surfaces. Its elliptic-bundle
structure prevents a direct identification with our generic
simple two-tori. For the related codimension-two extension
question, \citet[Theorem 7.1.3]{Nakayama2002Root} proves an
extension criterion for locally bimeromorphically projective
elliptic fibrations over a punctured polydisc: the fibration
must be bimeromorphically projective over that punctured base.
It rules out the sectionless elliptic example discussed there.
Our regular fibres have dimension two and no asserted elliptic
quotient, so that theorem does not establish a nonexistence
result for the present filling.
\end{remark}
